\documentclass[12pt,spanish,mexico]{article}
\usepackage[utf8]{inputenc}
\PassOptionsToPackage{hyphens}{url}\usepackage{hyperref}
\usepackage[es-nodecimaldot]{babel} % nombre de opción corregido
\usepackage{amsthm, amssymb, amsmath, amsfonts, bbm, mathtools}
\usepackage{csquotes}
\usepackage{enumerate}
\usepackage[shortlabels]{enumitem}
\usepackage{titling}
\usepackage{blindtext}
\usepackage{lmodern} 
\usepackage{float}
\usepackage{mathrsfs}
\usepackage[capitalise]{cleveref}
\usepackage{caption} 
\usepackage{graphicx}
\usepackage{chngcntr}
\counterwithin{figure}{section}
\numberwithin{equation}{section}
\usepackage[vmargin=1in, hmargin=0.5in]{geometry}  

\setlength{\topmargin}{-50pt}
\setlength{\headheight}{38.23112pt}
\setlength{\headsep}{24pt}

\makeatletter
\def\iftitle{\ifx\thetitle \empty \else }
\def\safetitle{\iftitle{\thetitle}\fi}
\def\headertitle{%
    \ifdefined\hw{\hw}\fi%
    \ifdefined\hw \iftitle{ -- }\fi%
    \safetitle%
}
\def\@maketitle{%
  \begin{center}%
Tarea 1 - Problemas inversos - 11/02/2026 \\
Muñoz Macías Christian Jaime 
    \vskip 1em%
    \hrule
  \end{center}%
  \par
  \vskip 1.5em%
}
\makeatother
\usepackage{amsmath}
\usepackage{tikz}

\newcommand\RADot{%
  \mathord{%
    \mspace{1mu}%
    \text{\radot}%
    \mspace{1mu}%
  }%
}

\newcommand\radot{%
    \tikz[line cap=round,x=1ex,y=1ex,line width=0.3pt]
    {\draw (0,1) |- (1,0) (0.55,0) arc(0:90:0.55); \fill (0.23,0.23) circle (0.05);}%
}

%% encabezado personalizado
\usepackage{fancyhdr}
%% estilo plain personalizado
\fancypagestyle{plain}{%
    \fancyhf{} % limpia encabezado y pie de página
    \renewcommand{\headrulewidth}{0pt}
    \setlength{\headheight}{0pt}
    \fancyfoot[R]{\thepage}
}
%% estilo para todas las páginas
\pagestyle{fancy}
\fancyhf{}
\fancyhead[C]{Maestría en Matemáticas Aplicadas}
\fancyfoot[R]{\thepage}

%%%%%%%%%%%%%%%%%%%%%%%%
%% DEFINICIONES ÚTILES %%
%%%%%%%%%%%%%%%%%%%%%%%%
%% definiciones matemáticas
\newcommand{\N}{\mathbb{N}}
\newcommand{\Z}{\mathbb{Z}}
\newcommand{\Ha}{\mathcal{H}}
\newcommand{\R}{\mathbb{R}}
\newcommand{\Q}{\mathbb{Q}}
\newcommand{\C}{\mathbb{C}}
\newcommand{\PP}{\mathbb{P}}

\newtheorem{proposition}{Proposición}
\newtheorem{theorem}{Teorema}
\newtheorem{lemma}{Lema}
\newtheorem{corollary}{Corolario}

\theoremstyle{definition}
\newtheorem{problem}{Problema}
\newtheorem{definition}{Definición}[section]
\newtheorem*{example}{Ejemplo}

\theoremstyle{remark}
\newtheorem*{observation}{Observación}
\newtheorem*{note}{Nota}

\newenvironment{solution}{\begin{proof}[Solución]}{\end{proof}}
\newenvironment{dem}{\begin{proof}[Demostración]}{\end{proof}}
\renewcommand{\qedsymbol}{$\blacksquare$}

%% comandos útiles
\newcommand{\eps}{\varepsilon}
\newcommand{\To}{\longrightarrow}
\newcommand{\Eu}{\varPhi}
\newcommand{\Mapsto}{\longmapsto}

%% delimitadores emparejados
\DeclarePairedDelimiter{\set}{\{}{\}}
\DeclarePairedDelimiter{\floor}\lfloor\rfloor
\DeclarePairedDelimiter{\ceil}\lceil\rceil
\DeclarePairedDelimiter{\abs}\lvert\rvert
\DeclarePairedDelimiter{\gen}\langle\rangle
\DeclareMathOperator{\sgn}{sgn}
\DeclareMathOperator{\ord}{ord}

%% operador
\newcommand{\op}[1]{\operatorname{#1}}

%% etiquetas del documento
\title{Tarea 1 - Problemas inversos - 11/02/2026}
\def\course{Materia}
\def\university{CIMAT - MsC. Maestría en matemáticas aplicadas}
\author{Muñoz Macías Christian Jaime}
\def\email{\url{christian.munoz@cimat.mx}}
\date{\today}

%% permitir saltos de línea en ecuaciones
\allowdisplaybreaks[4]

%%%%%%%%%%% INICIO DEL DOCUMENTO %%%%%%%%%%% 

\begin{document}
\maketitle

\textbf{1. Deconvolución.} Considera el problema de deconvolución presentado en la sección 1.1 del libro de Vogel

\[
g(x) = \int_{0}^{1} k(x - x') f(x') dx' \stackrel{\text{def}}{=} (Kf)(x), \quad 0 < x < 1 \quad (1)
\]

con núcleo Gaussiano

\[
k(x) = C \exp(-x^2/2\gamma^2)
\]

donde \(\gamma = 0.025\) y \(C = 1/\sqrt{\gamma^2 2\pi}\). Consideremos además la función

\[
f_{\text{true}}(x) = 
\begin{cases} 
0.75, & 0.1 < x < 0.25 \\
0.25, & 0.3 < x < 0.32 \\
\sin^4(2\pi x), & 0.5 < x < 1 \\
0, & \text{en otro caso}
\end{cases}
\]

Realiza la discretización del problema de deconvolución de acuerdo al libro de Vogel.

El problema inverso consiste en estimar \( f_{\text{true}} (x) \) dadas observaciones ruidosas de \( g \).

\begin{solution}

Sea $n \in \mathbb{N}$ y considera una partición uniforme del intervalo $(0,1)$ dada por

\[
x_j = \frac{j}{n}, \quad j=1,\dots,n,
\]

con paso de malla

\[
h = \frac{1}{n}.
\]

Aproximamos la integral en (1) mediante una suma de Riemann:

\[
g(x_i)
\approx
\sum_{j=1}^{n}
k(x_i - x_j)\, f(x_j)\, h,
\quad i=1,\dots,n.
\]

Definimos los vectores discretos

\[
f_j = f(x_j), \quad g_i = g(x_i),
\]

y la matriz

\[
K_{ij} = h\, k(x_i - x_j),
\quad i,j=1,\dots,n.
\]

el problema integral continuo se aproxima por el sistema lineal

\[
g = K f,
\]

donde $K \in \mathbb{R}^{n \times n}$ representa la discretización del operador integral $K$.

La función verdadera se discretiza como

\[
f_{\text{true},j}
=
f_{\text{true}}(x_j),
\quad j=1,\dots,n,
\]

y los datos se obtienen por

\[
g = K f_{\text{true}}.
\]


\end{solution}

\begin{enumerate}[label=(\alph*)]
    \item Lee la subsección 1.2.1 del libro y responde los ejercicios 1.5 y 1.6.
    \begin{enumerate}
        \item \textbf{Ejercicio 1.5.} Verifica que la ecuación (1.21) se satisface tanto para la función de filtro TSVD (1.11) como para la función de filtro de Tikhonov (1.13).
        \begin{equation}
        w_\alpha(s^2)s^{-1} \leq \alpha^{-1/2}
        \tag{1.21}
        \end{equation}
        \begin{equation}
        w_\alpha(s^2)=\begin{cases}
            1 & \text{si } s^2 > \alpha,\\
            0 & \text{si } s^2 \leq \alpha.
        \end{cases}
        \tag{1.11}
        \end{equation}
        \begin{equation}
            w_\alpha(s^2) =\frac{s^2}{s^2+\alpha}.
            \tag{1.13}
        \end{equation}
        \begin{solution}
            Para la función de filtro TSVD, si $(s^2 > \alpha)$ entonces
            $$w_\alpha(s^2)=1$$
            aplicando $s^{-1}$ en los dos lados de la ecuación anterior, obtenemos
            $$w_\alpha(s^2)s^{-1} =1 s^{-1}$$
            si $s^2 > \alpha$, entonces $s> \sqrt{\alpha}$ y por lo tanto $s^{-1} < \alpha^{-1/2}$, lo que implica que
            $$w_\alpha(s^2)s^{-1}< \alpha^{-1/2}$$
            lo que cumple la desigualdad (1.21).
            ahora si $s^2 < \alpha$ la función de filtro TSVD es
            $$w_\alpha(s^2)=0$$
            aplicando $s^{-1}$ en los dos lados de la ecuación anterior, obtenemos
            $$w_\alpha(s^2)s^{-1} =0 s^{-1}=0$$
            entonces $0<\alpha^{-1/2}$, esto cumple la desigualdad (1.21), por lo tanto la función de filtro TSVD cumple
            la desigualdad (1.21).\\
            Para la funcion de filtro de Tikhonov, tenemos que multiplicar ambos lados de la función  por $s^{-1}$.
            \begin{align*}
            w_\alpha(s^2)s^{-1} &=\frac{s^2}{s^2+\alpha}s^{-1}\\
            &=\frac{s}{s^2+\alpha}\\
            \frac{s}{s^2+\alpha} &\leq \alpha^{-1/2} \\
            s &\leq \alpha^{-1/2}(s^2+\alpha) \\
            0 &\leq \alpha^{-1/2}(s^2+\alpha)-s\\
            0 &\leq \alpha^{-1/2}s^2-s+\alpha^{1/2}\\
            0 &\leq s^2-\alpha^{1/2}s+\alpha\\
            0 &\leq s^2-\alpha^{1/2}s+\alpha+\frac{\alpha}{4}-\frac{\alpha}{4}\\
            0 &\leq \left(s^2-\frac{\alpha^{1/2}}{2}\right)^2+\frac{3 \alpha}{4}
            \end{align*}
            la última desigualdad es verdadera para todo $s$  por lo tanto cumple la desigualdad (1.21).
        \end{solution}
        \item \textbf{Ejercicio 1.6.} Usando (1.18) y (1.5), muestra que la ecuación (1.22) se cumple.
        \begin{equation}
            \begin{aligned}
                e_\alpha^{\text{ruido}} &= R_\alpha \eta\\
                &=\sum_{i=1}^{n}w_\alpha(s_i^2)s_i^{-1}(u_i^T \eta)v_i.
            \end{aligned}
            \tag{1.18}
        \end{equation}
        \begin{equation}
            \delta = \|\eta\|>0
            \tag{1.5}
        \end{equation}
        \begin{equation}
            \|e_\alpha^{\text{ruido}}\| \leq \alpha^{-1/2} \delta
            \tag{1.22}
        \end{equation}
        \begin{solution}
        De (1.18),

\[
e_\alpha^{\text{ruido}}
=
\sum_{i=1}^{n}
w_\alpha(s_i^2)s_i^{-1}(u_i^T \eta)v_i.
\]

Como $\{v_i\}$ es un sistema ortonormal,

\begin{align*}
\|e_\alpha^{\text{ruido}}\|^2
&=\left\langle \sum_{i=1}^{n}
w_\alpha(s_i^2)s_i^{-1}(u_i^T \eta)v_i,\sum_{i=1}^{n}
w_\alpha(s_i^2)s_i^{-1}(u_i^T \eta)v_i \right\rangle\\
&=\sum_{i=1}^{n}\left(w_\alpha(s_i^2)s_i^{-1}\right)^2 (u_i^T \eta)^2.
\end{align*}

Por el Ejercicio 1.5,

\[
\left|
w_\alpha(s_i^2)s_i^{-1}
\right|
\le
\alpha^{-1/2}.
\]

Por lo tanto,

\[
\|e_\alpha^{\text{ruido}}\|^2
\le
\alpha^{-1}
\sum_{i=1}^{n}
(u_i^T \eta)^2.
\]

Como $\{u_i\}$ también es ortonormal,

\[
\sum_{i=1}^{n}
(u_i^T \eta)^2
=
\|\eta\|^2.
\]

Usando (1.5),

\[
\|\eta\| = \delta,
\]

Se obtiene

\[
\|e_\alpha^{\text{ruido}}\|^2
\le
\alpha^{-1}\delta^2.
\]

Aplicando raíz cuadrada,

\[
\|e_\alpha^{\text{ruido}}\|
\le
\alpha^{-1/2}\delta.
\]

Esto prueba (1.22).
        \end{solution}
    \end{enumerate}
    
    \item Lee la subsección 1.2.2 del libro y responde los ejercicios 1.7, 1.8 y 1.9.
    \begin{enumerate}
        \item \textbf{Ejercicio 1.7.} Confirma la desigualdad (1.26).
        \begin{equation}
            \begin{aligned}
                \|e_\alpha^{\text{trunc}}\|^2&=\sum_{i=1}^{n}(w_\alpha(s_i^2)-1)^2s_i^2|z^{T}u_i|^2 \\
                &\leq \max_{1 \leq i \leq n}(w_\alpha(s_i^2)-1)^2 s_i^2 \|z\|^2 \\
                &\leq \alpha \|z\|^2
            \end{aligned}
            \tag{1.26}
        \end{equation}
        \begin{solution}
        
En la base ortonormal $\{u_i\}$,

\[
\|e_\alpha^{\text{trunc}}\|^2
=
\sum_{i=1}^{n}
(w_\alpha(s_i^2)-1)^2
s_i^2
|z^{T}u_i|^2.
\]

Sea

\[
\beta
=
\max_{1\le i\le n}
(w_\alpha(s_i^2)-1)^2 s_i^2.
\]

Entonces cada término de la suma está acotado por $\beta|z^T u_i|^2$, y por lo tanto

\[
\|e_\alpha^{\text{trunc}}\|^2
\le
\beta
\sum_{i=1}^{n}
|z^{T}u_i|^2.
\]

Como $\{u_i\}$ es una base ortonormal,

\[
\sum_{i=1}^{n}
|z^{T}u_i|^2
=
\|z\|^2.
\]

Por consiguiente,

\[
\|e_\alpha^{\text{trunc}}\|^2
\le
\max_{1\le i\le n}
(w_\alpha(s_i^2)-1)^2 s_i^2
\|z\|^2.
\]

Ahora verificamos que para filtros regulares TSVD se cumple

\[
(w_\alpha(s^2)-1)^2 s^2 \le \alpha
\quad \text{para todo } s.
\]


Si $s^2>\alpha$, entonces $w_\alpha(s^2)=1$ y

\[
(1-1)^2 s^2 = 0 \le \alpha.
\]

Si $s^2\le\alpha$, entonces $w_\alpha(s^2)=0$ y

\[
(0-1)^2 s^2 = s^2 \le \alpha.
\]

Por lo tanto, la desigualdad se cumple para todo $s$ en el caso TSVD.

Concluimos que

\[
(w_\alpha(s^2)-1)^2 s^2 \le \alpha
\quad \text{para todo } s
\]

tanto para TSVD.
Por lo tanto

\[
\|e_\alpha^{\text{trunc}}\|^2
\le
\alpha \|z\|^2.
\]

Esto confirma (1.26).

        \end{solution}
        \item \textbf{Ejercicio 1.8.} Muestra que la desigualdad (1.26) es óptima (sharp). Para esto, proporciona el vector \(z\) para el cual se alcanza la igualdad.
\begin{solution}

Sea

\[
\beta^*
=
\max_i
(w_\alpha(s_i^2)-1)^2 s_i^2
\]

y sea $i^*$ un índice donde se alcanza el máximo. Tomamos

\[
z = u_{i^*}.
\]

Como $\{u_i\}$ es ortonormal,

\[
z^T u_{i^*}=1,
\quad
z^T u_i=0 \ (i\neq i^*),
\quad
\|z\|=1.
\]

Entonces

\[
\|e_\alpha^{\text{trunc}}\|^2
=
(w_\alpha(s_{i^*}^2)-1)^2 s_{i^*}^2
=
\beta^*
=
\beta^{*}\|z\|^2.
\]

Por lo tanto la desigualdad (1.26) se alcanza como igualdad,
y la cota es óptima (sharp).   
\end{solution}

        \item \textbf{Ejercicio 1.9.} Muestra que el lado derecho de (1.27) se minimiza con la elección (1.28). Luego confirma (1.29).
        \begin{equation}
            \|e_\alpha\| \leq \alpha^{1/2}\|z\|+\alpha^{-1/2} \delta.
            \tag{1.27}
        \end{equation}
        \begin{equation}
            \alpha = \frac{\delta}{\|z\|}.
            \tag{1.28}
        \end{equation}
        \begin{equation}
            \|e_\alpha\|\leq 2\|z\|^{1/2} \delta^{1/2}
            \tag{1.29}
        \end{equation}
    \end{enumerate}
    \begin{solution}
Queremos minimizar para $\alpha>0$ la función

\[
f(\alpha)=\alpha^{1/2}\|z\|+\alpha^{-1/2}\delta.
\]

Derivando respecto a $\alpha$,

\[
f'(\alpha)
=
\frac{1}{2}\alpha^{-1/2}\|z\|
-
\frac{1}{2}\alpha^{-3/2}\delta.
\]

Igualamos a cero:

\begin{align*}
0
&=
\alpha^{-1/2}\|z\|
-
\alpha^{-3/2}\delta \\
0
&=
\alpha^{-1/2}\left(\|z\|-\alpha^{-1}\delta\right).
\end{align*}

Como $\alpha^{-1/2}\neq 0$, se obtiene

\[
\|z\|-\alpha^{-1}\delta=0
\quad\Rightarrow\quad
\alpha=\frac{\delta}{\|z\|}.
\]

Además,
\begin{align*}
f''(\alpha)
&=
-\frac{1}{4}\alpha^{-3/2}\|z\|+
\frac{3}{4}\alpha^{-5/2}\delta,\\
f''\left(\frac{\delta}{\|z\|}\right)
&=
-\frac{1}{4}\left(\frac{\delta}{\|z\|}\right)^{-3/2}\|z\|+
\frac{3}{4}\left(\frac{\delta}{\|z\|}\right)^{-5/2}\delta,\\
&=
-\frac{1}{4}\delta^{-3/2}\|z\|^{5/2}+
\frac{3}{4}\delta^{-3/2}\|z\|^{5/2},\\
&=
\frac{1}{2}\delta^{-3/2}\|z\|^{5/2},\\
&=
\frac{1}{2}\|z\|\left(\frac{\|z\|}{\delta}\right)^{3/2},\\
&=
\frac{1}{2}\|z\|\alpha^{-3/2}\frac{\delta}{\delta},\\
&=\frac{1}{2}\alpha^{-5/2}\delta,\\
\end{align*}
como $\alpha$ y $\delta$ son positivos, se concluye que 
\[
f''(\alpha)>0,
\]

por lo que este punto crítico corresponde a un mínimo.
Esto prueba (1.28).

Sustituyendo en (1.27):

\begin{align*}
\|e_\alpha\|
&\leq
\left(\frac{\delta}{\|z\|}\right)^{1/2}\|z\|
+
\left(\frac{\delta}{\|z\|}\right)^{-1/2}\delta \\
&\leq
\delta^{1/2}\|z\|^{1/2}
+
\delta^{1/2}\|z\|^{1/2} \\
&\leq
2\delta^{1/2}\|z\|^{1/2}.
\end{align*}

Esto confirma (1.29).
\end{solution}
    \item Haz un reporte sobre la solución del problema inverso con regularización TSVD y Tikhonov con principio de discrepancia, usando diferentes tamaños de discretización que tú elijas (por ejemplo \( n = 10, 50, 100 \) puntos y relación señal-ruido \( SNR = \frac{\max}{\sigma_g} = 10, 20, 100 \)).
\begin{figure}[H]
    \centering
    \includegraphics[width=0.75\linewidth]{C:/Users/chris/OneDrive/Escritorio/python/CIMAT/problemas inversos/Tarea_1/Tarea_1c_n10_SNR10.pdf}
    \caption{Comparación de filtros ($n=10$, $\mathrm{SNR}=10$)}
\end{figure}
\begin{figure}[H]
    \centering
    \includegraphics[width=0.75\linewidth]{C:/Users/chris/OneDrive/Escritorio/python/CIMAT/problemas inversos/Tarea_1/Tarea_1c_n10_SNR20.pdf}
    \caption{Comparación de filtros ($n=10$, $\mathrm{SNR}=20$)}
\end{figure}
\begin{figure}[H]
    \centering
    \includegraphics[width=0.75\linewidth]{C:/Users/chris/OneDrive/Escritorio/python/CIMAT/problemas inversos/Tarea_1/Tarea_1c_n10_SNR100.pdf}
    \caption{Comparación de filtros ($n=10$, $\mathrm{SNR}=100$)}
\end{figure}
\begin{figure}[H]
    \centering
    \includegraphics[width=0.75\linewidth]{C:/Users/chris/OneDrive/Escritorio/python/CIMAT/problemas inversos/Tarea_1/Tarea_1c_n50_SNR10.pdf}
    \caption{Comparación de filtros ($n=50$, $\mathrm{SNR}=10$)}
\end{figure}
\begin{figure}[H]
    \centering
    \includegraphics[width=0.75\linewidth]{C:/Users/chris/OneDrive/Escritorio/python/CIMAT/problemas inversos/Tarea_1/Tarea_1c_n50_SNR20.pdf}
    \caption{Comparación de filtros ($n=50$, $\mathrm{SNR}=20$)}
\end{figure}
\begin{figure}[H]
    \centering
    \includegraphics[width=0.75\linewidth]{C:/Users/chris/OneDrive/Escritorio/python/CIMAT/problemas inversos/Tarea_1/Tarea_1c_n50_SNR100.pdf}
    \caption{Comparación de filtros ($n=50$, $\mathrm{SNR}=100$)}
\end{figure}
\begin{figure}[H]
    \centering
    \includegraphics[width=0.75\linewidth]{C:/Users/chris/OneDrive/Escritorio/python/CIMAT/problemas inversos/Tarea_1/Tarea_1c_n100_SNR10.pdf}
    \caption{Comparación de filtros ($n=100$, $\mathrm{SNR}=10$)}
\end{figure}
\begin{figure}[H]
    \centering
    \includegraphics[width=0.75\linewidth]{C:/Users/chris/OneDrive/Escritorio/python/CIMAT/problemas inversos/Tarea_1/Tarea_1c_n100_SNR20.pdf}
    \caption{Comparación de filtros ($n=100$, $\mathrm{SNR}=20$)}
\end{figure}
\begin{figure}[H]
    \centering
    \includegraphics[width=0.75\linewidth]{C:/Users/chris/OneDrive/Escritorio/python/CIMAT/problemas inversos/Tarea_1/Tarea_1c_n100_SNR100.pdf}
    \caption{Comparación de filtros ($n=100$, $\mathrm{SNR}=100$)}
\end{figure}
\textbf{Conclusión}
Se observa que para la discretización  $n=10$ los valores singulares se mantienen por encima de 1 y no presentan un decaimiento apreciable en ninguna de las pruebas de SNR. Esto indica que la matriz obtenida con esta discretización es numéricamente estable y poco sensible al ruido.\\

En cambio, para $n=50$ y $n=100$ los valores singulares muestran un decaimiento marcado. A mayor SNR, el espectro se extiende hacia valores más pequeños, lo que implica que la matriz se vuelve más sensible a perturbaciones en los datos. Esta sensibilidad se traduce en una mayor amplificación del ruido cuando se intenta invertir el sistema sin regularización.

\end{enumerate}

\end{document}